\documentclass[10pt]{article}
\usepackage[T1]{fontenc}
\usepackage[utf8]{inputenc}
\usepackage[french]{babel}
\usepackage{amsmath,amssymb,mathrsfs,mathpazo}
\usepackage{geometry}
\geometry{a4paper,margin=2cm}
\usepackage{graphicx,array,tabularx,tabularray,colortbl,multirow,multicol,booktabs,enumitem,pifont,ulem,fancybox,diagbox,cancel}
\usepackage{pstricks,pstricks-add,pst-plot,pst-tree,pst-node,pst-text,pst-func,pst-math,pst-bspline,pst-3dplot,pst-eucl,pst-solides3d}
\let\pscalculate\undefined
\usepackage{pst-bezier}
\usepackage{tikz,pgfplots,tkz-tab}
\usetikzlibrary{arrows,arrows.meta,patterns,calc,positioning,trees,decorations.markings,shapes.geometric,shapes.misc,shapes.arrows,intersections,decorations.pathmorphing,decorations.pathreplacing,backgrounds,shadows,babel}
\pgfplotsset{compat=1.18}
\usepgfplotslibrary{fillbetween}
\usepackage[np]{numprint}
\usepackage{listings,makecell,pas-tableur,twoopt}
\let\barre\undefined
\usepackage[pstricks]{tablvar}
\definecolor{darkgreen}{RGB}{0,100,0}
\providecommand{\red}{\color{red}}
\providecommand{\blue}{\color{blue}}
\providecommand{\green}{\color{green}}
\providecommand{\black}{\color{black}}
\providecommand{\white}{\color{white}}
\providecommand{\gray}{\color{gray}}
\providecommand{\cyan}{\color{cyan}}
\providecommand{\magenta}{\color{magenta}}
\providecommand{\yellow}{\color{yellow}}
\providecommand{\violet}{\color{violet}}
\providecommand{\orange}{\color{orange}}
\providecommand{\euro}{\texteuro}
\providecommand{\EUR}{\texteuro}
\providecommand{\ieme}{\textsuperscript{e}}
\providecommand{\degres}{\ensuremath{^\circ}}
\providecommand{\degree}{\ensuremath{^\circ}}
\providecommand{\vv}[1]{\overrightarrow{#1}}
\providecommand{\Coord}[1]{#1}
\providecommand{\decosix}{}
\providecommand{\starredbullet}{\ensuremath{\bullet}}
\pagestyle{empty}
\setlength{\parindent}{0pt}
\usepackage{adjustbox}\DeclareUnicodeCharacter{2236}{:}\DeclareUnicodeCharacter{2190}{\ensuremath{\leftarrow}}
\providecommand{\R}{\mathbb{R}}
\providecommand{\N}{\mathbb{N}}
\providecommand{\Z}{\mathbb{Z}}
\providecommand{\Q}{\mathbb{Q}}
\providecommand{\D}{\mathbb{D}}
\providecommand{\C}{\mathbb{C}}
\providecommand{\e}{{\mathrm{e}}}
\providecommand{\vect}[1]{\overrightarrow{#1}}
\providecommand{\vectt}[1]{\overrightarrow{\mathrm{#1}}}
\providecommand{\vvt}[1]{\overrightarrow{\mathrm{#1}}}
\providecommand{\np}[2][]{{#2~#1}}
\providecommand{\vv}[1]{{\overrightarrow{#1}}}
\renewcommand{\d}{{\mathrm{d}}}
\providecommand{\degres}{{^\circ}}
\providecommand{\degre}{{^\circ}}
\providecommand{\blue}{\color{blue}}
\providecommand{\red}{\color{red}}
\providecommand{\enskip}{\quad}
\providecommand{\llbracket}{\lbrack\!\lbrack}
\providecommand{\rrbracket}{\rbrack\!\rbrack}
\providecommand{\euro}{{\text{€}}}
\providecommand{\no}{{\text{n°}}}
\providecommand{\interval}[2]{{[#1;#2]}}
\providecommand{\intervFO}[2]{{[#1;#2[}}
\providecommand{\dotfill}{\ldots}

\providecommand{\ds}{\displaystyle}
\providecommand{\textding}[1]{#1}
\providecommand{\Oij}{\(\mathrm{O};\vec{\imath},\vec{\jmath}\)}
\providecommand{\Oijk}{\(\mathrm{O};\vec{\imath},\vec{\jmath},\vec{k}\)}
\providecommand{\Ouv}{\(\mathrm{O};\vec{u},\vec{v}\)}
\providecommand{\intervFF}[2]{\ensuremath{[#1;#2]}}
\renewcommand{\euro}{\ifmmode\text{\texteuro}\else\texteuro\fi}
\DeclareUnicodeCharacter{00B0}{\ifmmode^{\circ}\else\textdegree\fi}
\DeclareUnicodeCharacter{20AC}{\ifmmode\text{\texteuro}\else\texteuro\fi}
% Énoncés : Éducation nationale. Transcription et corrigés : APMEP, Denis Vergès et contributeurs. Découpage et classement : CoopMaths, Sébastien Lozano, Éric Elter et contributeurs. Adaptation de la présentation : Maths Explorer.
% CC BY-SA : https://coopmaths.fr/www/about/
% https://coopmaths.fr/alea/static/bac/2021/tex/bac_2021_06_asie_10.tex
\begin{document}
\noindent\textit{APMEP / CoopMaths. CC BY-SA. Présentation adaptée par Maths Explorer.}\par\medskip
\medskip

\begin{tabular}{|l|}\hline
\textbf{Principaux domaines abordés}\\
-- Suites, étude de fonction\\
-- Fonction logarithme\\ \hline
\end{tabular}

\medskip

Soit la fonction $f$ définie sur l'intervalle $]1~;~ +\infty[$ par 

\[f(x) = x - \ln (x - 1).\]

On considère la suite $\left(u_n\right)$ de terme initial $u_0 = 10$ et telle que $u_{n+1} = f\left(u_n\right)$ pour tout entier naturel $n$.

\bigskip

\textbf{Partie I :}

\medskip

La feuille de calcul ci-dessous a permis d'obtenir des valeurs approchées des premiers termes de la suite $\left(u_n\right)$.

\begin{center}
\begin{tabularx}{0.6\linewidth}{|c|*{2}{>{\centering \arraybackslash}X|}}\hline
&A &B\\ \hline
1 &$n$&$u_n$\\ \hline
2 &0&10\\ \hline
3& 1&\np{7,80277542}\\ \hline
4& 2&\np{5,88544474}\\ \hline
5& 3&\np{4,29918442}\\ \hline
6& 4&\np{3,10550913}\\ \hline
7& 5&\np{2,36095182}\\ \hline
8& 6&\np{2,0527675}\\ \hline
9& 7&\np{2,00134509}\\ \hline
10& 8&\np{2,0000009}\\ \hline
\end{tabularx}
\end{center}

\medskip

\begin{enumerate}
\item Quelle formule a été saisie dans la cellule B3 pour permettre le calcul des valeurs approchées de $\left(u_n\right)$ par recopie vers le bas ?
\item À l'aide de ces valeurs, conjecturer le sens de variation et la limite de la suite $\left(u_n\right)$.
\end{enumerate}

\bigskip

\textbf{Partie II :}

\medskip

On rappelle que la fonction $f$ est définie sur l'intervalle $]1~;~ +\infty[$ par 

\[f(x) = x - \ln (x - 1).\]

\medskip

\begin{enumerate}
\item Calculer $\displaystyle\lim_{x \to 1} f(x)$. On admettra que $\displaystyle\lim_{x \to + \infty} f(x) = + \infty$.
\item  
	\begin{enumerate}
		\item Soit $f'$ la fonction dérivée de $f$. Montrer que pour tout $x \in ]1~;~ +\infty[$,\,  $f'(x) = \dfrac{x - 2}{x - 1}$.
		\item En déduire le tableau des variations de $f$ sur l'intervalle $]1~;~ +\infty[$, complété par les limites.
		\item Justifier que pour tout $x
\geqslant  2$,\,  $f(x) \geqslant  2$.
	\end{enumerate}
\end{enumerate}

\bigskip

\textbf{Partie III :}

\medskip

\begin{enumerate}
\item En utilisant les résultats de la partie II, démontrer par récurrence que $u_n \geqslant  2$ pour tout entier naturel $n$.
\item Montrer que la suite $\left(u_n\right)$ est décroissante.
\item En déduire que la suite $\left(u_n\right)$ est convergente. On note $\ell$ sa limite.
\item On admet que $\ell$ vérifie $f(\ell) = \ell$. Donner la valeur de $\ell$.
\end{enumerate}

\end{document}
